\chapter{第十二讲\quad 共同犯罪}
\label{lecture:12}
\noindent{\small\hyperref[ref:lecture:12]{查阅本讲扩展规则}}\par

先将案件写成事件表 \((\text{人},\text{行为},\text{时间},\text{对象},\text{罪过})\)，再逐人连线。连线须写明是共同实行、支配、引起犯意还是促进犯罪，以及其关联的行为和结果。连线成立后仍须独立检查各人的责任与量刑。以下按教材所用的违法事实层与个人责任层分别推导。法条层的共同犯罪须满足《刑法》第25条关于共同故意的要求。

\section{第一节\quad 共同犯罪的分类与原理}

\paragraph{法条和模型}
对至少两人的集合 \(G\)，令 \(W(G,g)\) 为各人以分工共同作用于罪 \(g\) 的违法事实，\(I(G,g)\) 为相应的共同犯罪故意。以 \(C(G,g)\) 为共同犯罪的基本判断，先检验
\[
|G|\ge2,\qquad W(G,g),\qquad I(G,g).
\]
普通共同犯罪须有客观共同参与和主观共同故意。教材把在违法层有联系与在责任层可处罚分开。作用上的主犯、从犯、胁从犯与分工上的正犯、教唆犯、帮助犯，分别记录。多人作为正犯共同实行的属简单型，正犯外另有教唆、帮助的属复杂型。犯罪集团还需第26条规定的组织条件。

\paragraph{推导}
甲在乙知情并共同计划的情况下压制被害人反抗，乙当场取财，二人均有非法占有目的。甲与乙的动作分别实现抢劫的强制、取财环节，合为 \(W(G,\text{抢劫})\)，共同计划和依约实行证成 \(I(G,\text{抢劫})\)，故在其他要件齐备时二人共同犯罪。若二人互不知情地在相邻店铺各自盗窃，相同的罪名不能推出 \(W\) 或 \(I\)，须按各自行为评价。若只是共同过失，第25条第2款要求分别判断各自所犯罪名。

\section{第二节\quad 共同正犯}

\paragraph{教材模型}
令 \(J(x,y,g)\) 为二人在共同计划内分担具有实行意义的任务，\(B(g)\) 为实行造成相应结果，\(F(x,g)\) 为该结果在 \(x\) 的共同实行与可归责范围内。共同正犯归责的核心是
\[
J(x,y,g)\land B(g)\land F(x,g)
\Longrightarrow \text{可对 }x\text{ 归属该结果}.
\]
“部分实行、全部负责”的全部范围由共同实行和可归责性界定。加重结果须另查每人的相应罪过及法律要求。

\paragraph{推导}
甲、乙有共同伤害故意，并在同一围打中分工控制、殴打丙，均对重伤有相应罪过。重伤来自这一共同攻击，证据虽不能确定致伤动作具体出自哪一人，但已证明处于共同实行范围。于是 \(J\land B\land F(\text{甲},g)\land F(\text{乙},g)\) 成立，分别适用规则归责。若二人并无联系，均以为瞄准猎物，各自射击而无法查明哪一枪致人死亡，不能使用共同故意的规则补证个人归责。

\paragraph{争议的分叉}
若二人约好一起射击但都误把人当猎物，必须承认第25条第2款的界限。教材展示三种模型。其一要求共同故意，且认为无法单独证明死亡归属，故在此假设下不能定过失致人死亡。其二把共同过失纳入共同正犯，存在与第25条第2款不合的问题。其三区分共同行为的结果归责与共同犯罪的标签，先论各人的共同风险贡献，再分别判断各自过失犯罪。答题注明采用何说，并写出能否证明各人风险贡献的理由，不应把有联络直接等同于有共同犯罪故意。

\section{第三节\quad 间接正犯}

\paragraph{教材模型}
令 \(D(x,y,b)\) 为幕后者 \(x\) 对 \(y\) 实施行为 \(b\) 有实行支配，\(K(x,D)\) 为 \(x\) 明知且利用这种支配，\(H(x,g)\) 为 \(x\) 满足所论罪名的正犯资格。教材典型模型为
\[
D(x,y,b)\land K(x,D)\land H(x,g)\land
\text{\(b\) 实现相应实行行为}
\Longrightarrow \text{\(x\) 可作为间接正犯评价}.
\]
支配可由强制、利用事实错误等产生。身份犯另查主体资格及具体身份规范。教材也讨论有身份者利用无身份者实施身份犯罪的学理路径，不能由任何身份差异直接推出间接正犯。

\paragraph{推导}
甲故意把毒物包装成普通药品，托完全不知情的乙给丙服用，丙服后因毒物死亡。乙对关键事实的错误由甲制造，甲明知并利用错误控制给药过程，故 \(D\land K\) 成立。若 \(H\) 以及杀人故意、实行、死亡归责等均成立，甲可作为故意杀人间接正犯既遂。若乙知道全部事实且独立决定杀人，仍须证明甲对实行有额外支配。只证明甲引起乙的犯意，通常转查教唆，不能把有指使直接换成有支配。

\section{第四节\quad 正犯与共犯的关系}

\paragraph{教材模型}
狭义共犯指教唆犯、帮助犯。令 \(P(y,g)\) 表示正犯侧至少实施了该罪相关的可评价违法行为，\(L(x,y,g)\) 表示教唆或帮助与该违法行为有法定关联。按教材的从属性模型，
\[
P(y,g)\land L(x,y,g)\land\text{\(x\) 有相应故意}
\Longrightarrow \text{狭义共犯的基础成立}.
\]
正犯侧是相关违法事实的来源。正犯个人是否达到年龄、是否可以谴责，与该事实的存在是两项判断。正犯的阶段影响共犯可评阶段，正犯的中止或其他人身性事由须分人审查。教唆还需结合第29条及其解释分歧。

\paragraph{推导}
乙为了杀人向甲索要工具，甲明知用途而提供，乙尚在准备阶段便自动销毁工具并放弃。按教材模型，乙有可评价的准备行为，甲的帮助作用到达这一阶段，但尚未发生着手。乙成立预备阶段的中止。甲没有退出，是因乙停止而无法继续，故按甲的作用与形态认定预备。若乙只是佯装接下提议、从未实施任何可评价违法行为，依一般从属性模型不具 \(P\)，结论不同。不能只凭教唆或帮助者的想法补足正犯事实。

\section{第五节\quad 教唆犯}

\paragraph{法条和模型}
《刑法》第29条。令 \(M(x,y,g)\) 为 \(x\) 的行为引起 \(y\) 实施罪 \(g\) 的犯意，\(L_g\) 为该犯意实际连接到 \(y\) 的相关违法行为，\(I_x\) 为 \(x\) 故意引起这项犯罪。普通教唆的教材模型要求 \(M\land L_g\land I_x\)。教唆已立与既遂是两步，后者另查由这一实行过程引起的结果。教唆未满18周岁者犯罪，应当从重处罚。刑罚另按教唆者在共同犯罪中的作用判断。

\paragraph{推导}
甲说服原本无盗窃犯意的乙窃取一件指定财物，乙由此接受计划并着手，后因被发现而失败。甲的劝说引起乙的具体犯意，给出 \(M\)，乙因此着手给出 \(L_g\)，已证故意给出 \(I_x\)。故可论盗窃的教唆，按教材从属性模型评价为未遂。若乙原本已坚定决定以相同方案盗窃，甲没有引起其犯意，不能推出 \(M\)。此时仅在证成真实促进作用时才转查心理性帮助。

\paragraph{第29条第2款的分歧}
被教唆的人没有犯被教唆之罪时，对教唆犯依法可以从轻或者减轻处罚。教材列出独立性与从属性的解释分歧。从属性解释强调正犯一侧的违法行为，教材把该款主要理解为已着手而未既遂等情形。独立性解释容许在一定范围内单独评价教唆行为。若题干是明确拒绝、后来另起犯意实施同样罪行，应先否定前后因果连接，再说明采取哪种条文解释，不得把后次结果直接归给早先劝说者。

\section{第六节\quad 帮助犯}

\paragraph{教材模型}
令 \(A(x,y,g)\) 为实际提供物理或心理促进，\(I(x,y,g)\) 为明知并故意促进相应犯罪，\(P(y,g)\) 为正犯的相关违法行为。先检查 \(A\land I\land P\)，再检查促进是否延续到实行及结果。促进可表现为降低作案障碍、降低被发现概率或使正犯因得到承诺而安心。不能把场所邻近、职业交往或口头表示支持本身直接等同于 \(A\)。

\paragraph{推导}
甲已知乙准备盗窃，依共同约定把能开启仓库的钥匙交给乙。乙借此开门并取走财物。钥匙在实际实行链中消除门锁障碍，给出 \(A\)，事前约定和对用途的认识给出 \(I\)，乙利用钥匙完成盗窃给出 \(P\) 与结果关联，甲可作为盗窃帮助犯既遂。若乙未用钥匙、根本不知甲准备支援，且没有任何其他促进作用，不能仅从乙既遂推出甲也既遂。需另审先前准备、作用消除及阶段。

\paragraph{中立行为边界}
普通运输、销售或技术服务要结合行为人所知、服务是否定向嵌入犯罪和具体促进来评价。例如按公开价格向陌生人卖普通手套且不知犯罪计划，不能推出 \(I\)。若特别罪名有独立构成要件，另按该条审查，不能以一般帮助犯的标签取代其要件。

\section{第七节\quad 正犯的责任年龄和责任能力}

\paragraph{教材模型}
令 \(U(y,g)\) 为 \(y\) 实施符合罪 \(g\) 的违法行为，\(R(y,g)\) 为其具备相应责任条件，\(L(x,y,g)\) 为 \(x\) 对该违法行为有可归责的参与，\(D(x,y)\) 为 \(x\) 支配 \(y\)。在教材分层模型中，
\[
U(y,g)\land\neg R(y,g)
\ \nRightarrow\ \neg L(x,y,g),\qquad
\neg R(y,g)\ \nRightarrow\ D(x,y).
\]
年龄不足或失能只回答实行者的个人责任。其他人的故意、贡献及对实行者有无支配，须分别证明。

\paragraph{推导}
乙未达盗窃的责任年龄，但能理解取财行为且独立选择作案，成年甲故意为乙提供能实际开门的钥匙。非法取财事实使 \(U\) 成立，年龄不足使 \(\neg R\) 成立，提供工具并实际促进则使 \(L\) 成立。按教材模型，可在甲的自身条件均满足时认定甲的帮助责任，乙不因本罪受刑事处罚。若进一步证明甲故意利用乙对财物权属的关键错误控制行动，则重新检验 \(D\) 与利用故意。不能只凭乙年龄小就将甲的帮助提升为间接正犯。

\section{第八节\quad 共同犯罪的参与时间}

\paragraph{教材模型}
设 \(t_x\) 为 \(x\) 介入的时点。结果 \(r\) 能否归给 \(x\)，不能只看 \(r\) 在不在 \(t_x\) 之后发生，还应看导致 \(r\) 的行为或风险是否在参与范围。若 \(x\) 事前无通谋，只能从介入后实际承接、参与或促进的部分寻找归责根据。退出时令 \(P_x,M_x\) 分别为继续存在的物理、心理贡献。教材要求检查
\[
\neg P_x\land\neg M_x
\]
是否成立，并依犯罪阶段、停止的自动性和防止结果的条件进一步判断中止。

\paragraph{推导}
甲已经非法入户，乙其后知情并答应在门外望风，甲在屋内实施抢劫。乙的望风给出介入后的帮助，却没有参与或促进此前的入户。故按教材模型，对乙审查基本抢劫的帮助责任，不能把既往入户这一加重事实自动归给乙。又如乙在原约定中向甲提供钥匙，后来取回原件却未能收回已配的副本，甲仍用副本作案，物理贡献 \(P_x\) 未消除。只说自己退出，不能推出中止。仅离开望风点而让正犯仍确信有人望风，也须检查 \(M_x\)。

\paragraph{中止有效性的两种判断}
区分共同犯罪事实已经既遂 \(K_g\) 与退出者本人承担既遂归责 \(K_x\)。教材在积极阻止而同伙也未得逞的特定情形中，容许不要求未得逞完全是退出者的功劳。教材还将共同压制后积极阻止、却被同伙打晕而同伙既遂的例题，评价为退出者个人中止。后一个模型允许 \(K_g\) 成立而依教材的归责处理另判 \(K_x\)，其结论属于解释选择。若采用要求消除先前贡献并实际防止结果的模型，仍须证明这些条件，不能从积极反对单独推出中止。

\section{第九节\quad 片面的共同犯罪}

\paragraph{争议模型}
\(K_{xy}\) 表示 \(x\) 知道自己参与 \(y\) 的犯罪，\(K_{yx}\) 表示 \(y\) 知道 \(x\) 的参与。片面情形有 \(K_{xy}\land\neg K_{yx}\)。片面帮助肯定说在正犯相应违法事实已证明，且有参与故意、贡献及归责范围的条件下，可承认 \(x\) 的单向参与。要求双向联络的否定说在缺少 \(K_{yx}\) 时，不承认相应的共同性。教材还区分片面帮助、片面教唆与片面实行，提醒有的观点对前两类与片面实行采取不同立场。

\paragraph{推导}
乙实行盗窃时，甲知情并暗中解除门禁锁定，使乙能够进入，乙不知甲参与。解除门禁降低实行障碍，\(K_{xy}\land\neg K_{yx}\) 与促进均成立。在题设已经给出乙的盗窃违法事实时，肯定说可评价甲为片面帮助犯，乙仍按自己的单独实行评价。双向联络说不因甲知情即认定共同犯罪。若甲仅在远处观望，既未实际介入，乙也不知道有人支持，连肯定说所需的促进都没有证成。题目如改变为甲暗中以暴力压制反抗、乙不知情而取财，应另审片面实行的肯定或否定分支，不能把帮助结论直接升级为共同正犯。

\section{第十节\quad 共同犯罪的结合问题}

\paragraph{四个筛选器}
令 \(U_x\) 为 \(x\) 实际实行或促进的事实集合，\(I_x\) 为其故意涵盖的事实集合。可候选归责范围先取 \(U_x\cap I_x\)，再审各罪的身份 \(H\)、因果归责和其他法定条件。这个交集仅是检验共同故意范围的记号，结果加重部分另查法定罪过。
\begin{enumerate}
\item 真正身份犯先查实行所需身份，再查无身份者能否以帮助或教唆参与。身份只影响处罚的情形，另查各人的加减条件。贪污与职务侵占交叉还须看《刑法》第382条第3款及相关司法解释。
\item 不作为参与须查特定义务、履行能力、实际促进和相应故意。只目睹他人犯罪不能推出帮助。
\item 实行过限应对超出共同计划的部分另行审查。事后知道、当场知道但没有接受或促进，均不能仅凭知情推出共同故意。实行不足时，只有重轻行为之间存在可评价的构成包含及相应故意时，才可能按重轻交集归责。
\item 对象错误或打击错误要分别以每人的认识与实际作用为输入，不能把实行者的认识自动赋给教唆者。
\end{enumerate}

\paragraph{推导}
甲故意给乙望风盗窃，乙独自另起犯意实施暴力取财。已证甲不知也未接受暴力，故暴力不在 \(I_{\text{甲}}\) 中，无法归给甲作为共同抢劫故意。若在教材的包容评价下仍有盗窃的共同取财事实和帮助关系，可以就该范围评价甲的盗窃帮助责任。若甲发现乙施暴后明确同意并实际协助压制反抗，则出现新的参与时点，须重算 \(U_{\text{甲}}\cap I_{\text{甲}}\)。同样，若甲曾教唆乙入户抢劫而乙只实施普通抢劫，基本抢劫在原教唆范围内，不以未实施的入户情节评价甲。

\section{第十一节\quad 共同犯罪的处罚规定}

\paragraph{法条和模型}
《刑法》第26至29条、第97条。令 \(G\) 为由三人以上为共同实施犯罪而组成的较固定组织，\(F_G\) 为能归属该集团的犯罪集合，\(P_x,O_x,D_x\) 为 \(x\) 参与、组织、指挥的犯罪集合。对组织、领导集团的首要分子，处罚范围为 \(F_G\)。对其他主犯，范围为 \(P_x\cup O_x\cup D_x\)。作用次要或辅助的从犯，应当从轻、减轻或者免除。依法认定为被胁迫参加犯罪的，按犯罪情节减轻或者免除。第29条同时规定教唆的按作用处罚及教唆未成年人犯罪的从重规则。

\paragraph{推导}
甲组织、领导一个三人以上的较固定犯罪集团，集团实施盗窃甲案和乙案。集团成员丙另外私下实施与集团无关的伤害案。由组织结构及甲的角色，甲作为集团首要分子属于主犯。\(F_G=\{\text{盗窃甲案},\text{盗窃乙案}\}\)，故按该集团两案处罚。单凭丙是成员，不能把其独立伤害案加到 \(F_G\)。若乙在集团两案中仅辅助其中一案，先核对乙可归责范围，再依第27条决定从宽。

\paragraph{边界}
分工名与作用名不在同一维度。策动并在共同犯罪中起主要作用的教唆者可为主犯。聚众犯罪的首要分子可能只有一人，需先查条文的处罚主体及是否成立共同犯罪。身体动作完全被强制、行为不能归属于本人时，须先解决行为与正当化问题，不直接认定为胁从犯。

\section*{怎么记与闭卷自测}
\textbf{先人时，后故意，再连线，查范围，分责任，定作用。}
在草稿纸上给每个人一行，画三列“加入前、参与中、退出后”。在每个行为旁写下他当时知道什么。给正犯、教唆、帮助分别标“实行或支配、引起、促进”。再圈出不能归给这个人的范围，例如加入前的事实、超出故意的暴力、已经消除的贡献。最后把个人责任与主从作用各写一列。
\begin{enumerate}
\item 甲教唆乙盗窃，乙未达该罪责任年龄，却独立决定、实际完成盗窃。能否一律把甲判为间接正犯？\par
答：不能。先查乙的违法事实与甲的教唆因果，另证甲有实行支配及利用故意。乙的年龄只直接决定乙的责任门。
\item 甲给乙钥匙后自愿退出，乙已经复制钥匙并用副本完成盗窃。甲退出声明能否当然排除对结果的责任？\par
答：不能。副本继续发挥原帮助的物理作用，要核查贡献有无真正消除，不能只凭声明认定中止。
\item 甲乙只共同商量盗窃，乙突然持刀伤人，甲完全不知。怎样分两层列结论？\par
答：先对共同取财范围审查共同违法和参与。再单列暴力作为乙的过限事实，不能补出甲的共同故意。乙及甲的最终罪名各按实际行为、故意与归责范围确定。
\end{enumerate}
